\documentclass{article}
\usepackage{amsmath, amssymb, amsthm}
\usepackage{hyperref}
\usepackage{geometry}
\geometry{margin=1in}

\title{Erd\H{o}s Problem \#750}
\date{Last updated: 2025-08-31}
\author{Source: erdosproblems.com}

\begin{document}
\maketitle

\noindent\textbf{Status:} OPEN \quad \textbf{No prize}

\noindent\textbf{Tags:} graph theory, chromatic number

\medskip

\noindent\textbf{Problem Statement:}

Let $f(m)$ be some function such that $f(m)\to \infty$ as $m\to \infty$. Does there exist a graph $G$ of infinite chromatic number such that every subgraph on $m$ vertices contains an independent set of size at least $\frac{m}{2}-f(m)$?




In \cite{Er69b} Erd\H{o}s conjectures this for $f(m)=\epsilon m$ for any fixed $\epsilon&gt;0$. This follows from a result of Erd\H{o}s, Hajnal, and Szemer\'{e}di \cite{EHS82}, as described by msellke in the comments. In \cite{ErHa67b} Erd\H{o}s and Hajnal prove this for $f(m)\geq cm$ for all $c&gt;1/4$. See also [75] .




References

[EHS82] Erd\H{o}s, P. and Hajnal, A. and Szemer\'{e}di, E., On almost bipartite large chromatic graphs . Theory and practice of combinatorics (1982), 117-123.

[Er69b] Erd\H{o}s, P., Problems and results in chromatic graph theory . Proof Techniques in Graph Theory (Proc. Second Ann Arbor Graph Theory Conf., Ann Arbor, Mich., 1968) (1969), 27-35.

[ErHa67b] Erd\H{o}s, P. and Hajnal, Andr\'{a}s, On chromatic graphs . Mat. Lapok (1967), 1--4.


\bigskip
\noindent\small{Source: \url{https://www.erdosproblems.com/750}}
\end{document}
